\subsection{Singularities Preview}

\begin{definition} \label{def:singular-local-ring-and-embedding-excess}
\hlt{Singular Local Rings and Embedding-Dimension Excess}\\
Let $R$ be a nonzero Noetherian commutative local ring with identity.
\begin{enumerate}[label=(\roman*)]
\setlength{\itemsep}{0pt}
\item The ring $R$ is \hlt{singular} if it is not regular. Equivalently, $\dim R\neq\operatorname{edim}(R)$ by Definition \ref{def:krull-dimension-and-regular-local-ring}(ii).
\item Assume $\dim R\in\N_0$ and $\dim R\leq\operatorname{edim}(R)$. The \hlt{embedding-dimension excess} is $\operatorname{exc}(R)\defeq\operatorname{edim}(R)-\dim R$.
\end{enumerate}
\end{definition}

\begin{proposition} \label{prop:positive-embedding-excess-detects-singularity}
\hlt{Positive Embedding-Dimension Excess Detects Singularity}\\
Let $R$ be a nonzero Noetherian commutative local ring with identity.\\
If $\dim R\in\N_0$ and $\dim R\leq\operatorname{edim}(R)$, then $R$ is singular $\Longleftrightarrow\operatorname{exc}(R)>0$.
\begin{proof}
\begin{enumerate}[label=\roman*.]
\setlength{\itemsep}{0pt}
\item[$(\Rightarrow)$] If $R$ is singular, then $\dim R\neq\operatorname{edim}(R)$ by Definition \ref{def:singular-local-ring-and-embedding-excess}(i). Thus $\dim R<\operatorname{edim}(R)$ by the dimension-order hypothesis.\\
The inequality $\operatorname{exc}(R)>0$ holds by Definition \ref{def:singular-local-ring-and-embedding-excess}(ii).
\item[$(\Leftarrow)$] If $\operatorname{exc}(R)>0$, then $\dim R<\operatorname{edim}(R)$ by Definition \ref{def:singular-local-ring-and-embedding-excess}(ii). Hence $R$ is singular by Definition \ref{def:singular-local-ring-and-embedding-excess}(i).
\end{enumerate}
\hlt{Structural Commentary}: The dimension order converts failure of regularity equality into positivity of the excess.\\
\hlt{Key Insight}: The excess vanishes exactly when Krull and embedding dimensions coincide.
\end{proof}
\end{proposition}

\begin{example} \label{ex:square-zero-singular-local-ring} \hlt{A Square-Zero Singular Local Ring}\\
Let $k$ be a field and set $R\defeq k[x,y]/\langle x^2,xy,y^2\rangle$. Write $\bar{x},\bar{y}\in R$ for the residue classes of $x,y$, and set $\mathfrak{m}\defeq\langle\bar{x},\bar{y}\rangle$.\\
Relation $(\forall\,r\in R)(\exists!\,(a,b,c)\in k^3)\,r=a+b\bar{x}+c\bar{y}$ holds by polynomial reduction. Equality $\mathfrak{m}^2=\{0_R\}$ holds by quotient relations.\\
For $(a,u)\in(k\setminus\{0_k\})\times\mathfrak{m}$, $(a+u)(a^{-1}-a^{-2}u)=1_R$. Thus $(\forall\,r\in R\setminus\mathfrak{m})(\exists\,s\in R)\,rs=1_R$.\\
The quotient $R/\mathfrak{m}\cong k$ is a field, so $\mathfrak{m}$ is maximal by Theorem \ref{thm:quotient-by-maximal-ideal-is-field}.\\
The implication $(\forall\,\mathcal{J}\trianglelefteq R)\,[\mathcal{J}\neq R\Rightarrow\mathcal{J}\subseteq\mathfrak{m}]$ holds because a proper ideal contains no invertible element.\\
Hence $R$ is local with maximal ideal $\mathfrak{m}$. The equality $\dim_kR=3$ holds by the unique normal form. Fix $\mathcal{J}\trianglelefteq R$.\\
Since $\mathcal{J}$ is a $k$-subspace of $R$, choose a finite $k$-basis by Corollary \ref{cor:finite-dimensional-spaces-and-subspaces-have-bases}(i)--(ii). This basis generates $\mathcal{J}$ as an $R$-module.\\
Hence $R$ is Noetherian by Theorem \ref{thm:noetherian-iff-submodules-finitely-generated}(ii), because $\mathcal{J}$ was arbitrary.\\
The inclusion $(\forall\,\mathfrak{p}\in\operatorname{Spec}R)\,\{\bar{x},\bar{y}\}\subseteq\mathfrak{p}$ holds by $\bar{x}^2=\bar{y}^2=0_R$ and primeness. Thus $\operatorname{Spec}R=\{\mathfrak{m}\}$.\\
The equality $\dim R=0$ holds by Definition \ref{def:krull-dimension-and-regular-local-ring}(i). The tuple $(\bar{x}+\mathfrak{m}^2,\bar{y}+\mathfrak{m}^2)$ is a $k$-basis of $\mathfrak{m}/\mathfrak{m}^2$.\\
Thus $\operatorname{edim}(R)=2$ by Definition \ref{def:embedding-dimension}. Therefore $\operatorname{exc}(R)=2>0$, and $R$ is singular by Proposition \ref{prop:positive-embedding-excess-detects-singularity}.
\end{example}
